% 竖版扉页：外观在此定义；书名、作者、英文书名与配图路径由入口提供。
% 由 layout-portrait.tex 末尾加载，沿用其字体与颜色；路径均相对于项目根目录。
% 绝对定位依赖 TikZ 的 remember picture，至少编译两轮（根目录 run.sh 已保证）。

\definecolor{PortraitTitlePageInk}{HTML}{1D1D1F}
\definecolor{PortraitTitlePageSecondary}{HTML}{6E6E73}
\definecolor{PortraitTitlePageHairline}{HTML}{D1D1D6}
% 唯一强调色与代码高亮的蓝灰一致。
\colorlet{PortraitTitlePageAccent}{PortraitCodeBlue}
\newfontfamily\PortraitTitlePageMediumLatin{PingFang SC Medium}
\newCJKfontfamily\PortraitTitlePageMediumCJK{PingFang SC Medium}
% 提交信息的眉题：小号大写，加宽字距。
\newfontfamily\PortraitTitlePageEyebrow{PingFang SC Medium}[LetterSpace=14]
% 提交说明原样输出：此字体族不启用 TeX 连字，-- 等不会被替换。
\newfontfamily\PortraitTitlePagePlainLatin{PingFang SC}

% 配图区域：等比缩放后居中；标题区与提交信息的位置不随图片比例变化。
\newlength{\PortraitTitlePageImageWidth}
\newlength{\PortraitTitlePageImageHeight}
\setlength{\PortraitTitlePageImageWidth}{130mm}
\setlength{\PortraitTitlePageImageHeight}{110mm}

\makeatletter
% 元信息：英文书名与配图路径；不写或参数为空时省略对应内容。
\newcommand*{\EnglishTitle}[1]{\gdef\PortraitTitlePage@EnglishTitle{#1}}
\newcommand*{\TitlePageImage}[1]{\gdef\PortraitTitlePage@Image{#1}}
\let\PortraitTitlePage@EnglishTitle\@empty
\let\PortraitTitlePage@Image\@empty
% PDF 文档属性同步入口的书名与作者；入口未设置时不写入。
\let\PortraitTitlePage@DefaultTitle\@title
\let\PortraitTitlePage@DefaultAuthor\@author
\AtBeginDocument{%
  \ifx\@title\PortraitTitlePage@DefaultTitle\else\hypersetup{pdftitle={\@title}}\fi
  \ifx\@author\PortraitTitlePage@DefaultAuthor\else\hypersetup{pdfauthor={\@author}}\fi}
\makeatother

% 提交信息文件由 scripts/commit-info.sh 经根目录 run.sh 写入编译输出目录，
% 四行依次为日期、短哈希、提交说明、clean 或 dirty。按原样逐行读取，
% 不经 TeX 转义；文件不存在或缺少日期、哈希时不显示提交信息。
\ExplSyntaxOn
\ior_new:N \g__portrait_titlepage_ior
\str_new:N \g__portrait_titlepage_date_str
\str_new:N \g__portrait_titlepage_hash_str
\str_new:N \g__portrait_titlepage_subject_str
\str_new:N \g__portrait_titlepage_state_str
\bool_new:N \g__portrait_titlepage_commit_bool
\cs_new_protected:Npn \__portrait_titlepage_get:N #1
  {
    \ior_if_eof:NTF \g__portrait_titlepage_ior
      { \str_gclear:N #1 }
      {
        \ior_str_get:NN \g__portrait_titlepage_ior \l_tmpa_str
        \str_gset_eq:NN #1 \l_tmpa_str
      }
  }
\NewDocumentCommand \PortraitTitlePageReadCommit { m }
  {
    \bool_gset_false:N \g__portrait_titlepage_commit_bool
    \file_if_exist:nT {#1}
      {
        \ior_open:Nn \g__portrait_titlepage_ior {#1}
        \__portrait_titlepage_get:N \g__portrait_titlepage_date_str
        \__portrait_titlepage_get:N \g__portrait_titlepage_hash_str
        \__portrait_titlepage_get:N \g__portrait_titlepage_subject_str
        \__portrait_titlepage_get:N \g__portrait_titlepage_state_str
        \ior_close:N \g__portrait_titlepage_ior
        \bool_lazy_or:nnF
          { \str_if_empty_p:N \g__portrait_titlepage_date_str }
          { \str_if_empty_p:N \g__portrait_titlepage_hash_str }
          { \bool_gset_true:N \g__portrait_titlepage_commit_bool }
      }
  }
% 默认读取本次编译的 <jobname>.commit；样张可在 \maketitle 前改读其他文件。
\AtBeginDocument { \PortraitTitlePageReadCommit { \c_sys_jobname_str .commit } }
\NewDocumentCommand \PortraitTitlePageIfCommitTF { +m +m }
  { \bool_if:NTF \g__portrait_titlepage_commit_bool {#1} {#2} }
% 眉题式三行：眉题“LAST COMMIT”；日期与哈希，dirty 时附圆角细框标签；提交说明。
% 不用分隔符串联，依靠换行、留白及字体差异区分标签与数值。
\NewDocumentCommand \PortraitTitlePageCommitText { }
  {
    { \PortraitTitlePageEyebrow \fontsize { 6.5 } { 8 } \selectfont LAST ~ COMMIT } \\ [ 1.5pt ]
    { \ttfamily \str_use:N \g__portrait_titlepage_date_str }
    \hspace { 0.8em }
    { \ttfamily \color { PortraitTitlePageAccent } \str_use:N \g__portrait_titlepage_hash_str }
    \str_if_eq:VnT \g__portrait_titlepage_state_str { dirty }
      { \hspace { 0.6em } \usebox \PortraitTitlePageDirtyBox }
    \str_if_empty:NF \g__portrait_titlepage_subject_str
      { \\ \PortraitTitlePagePlainLatin \str_use:N \g__portrait_titlepage_subject_str }
  }
\ExplSyntaxOff
% 未提交修改标签：0.4 pt 浅灰圆角细框，与代码卡片同一视觉语言。
% 在扉页的 tikzpicture 之外预先排入盒子，避免嵌套 TikZ。
\newsavebox{\PortraitTitlePageDirtyBox}
\newcommand*{\PortraitTitlePageMakeDirtyBox}{%
  \sbox{\PortraitTitlePageDirtyBox}{\tikz[baseline=(tag.base)]
    \node[draw=PortraitTitlePageHairline,line width=0.4pt,rounded corners=1.1mm,
      inner xsep=1.3mm,inner ysep=0.45mm,text=PortraitTitlePageInk,
      font=\sffamily\fontsize{6.5}{8}\selectfont] (tag) {含未提交修改};}}

% 版面以纸张中心为原点：标题区在上，配图居中偏下，提交信息贴近底边。
\makeatletter
\renewcommand{\maketitle}{%
  \clearpage
  \thispagestyle{empty}%
  \PortraitTitlePageMakeDirtyBox
  \null
  \begin{tikzpicture}[remember picture,overlay]
    \begin{scope}[shift={(current page.center)}]
      \node[anchor=base,inner sep=0pt,text=PortraitTitlePageInk,
        font=\sffamily\bfseries\fontsize{44}{52}\selectfont] at (0,80mm) {\@title};
      \ifx\PortraitTitlePage@EnglishTitle\@empty\else
        \node[anchor=base,inner sep=0pt,text=PortraitTitlePageSecondary,
          font=\sffamily\fontsize{15}{18}\selectfont] at (0,66.5mm)
          {\PortraitTitlePage@EnglishTitle};
      \fi
      \node[anchor=base,inner sep=0pt,text=PortraitTitlePageInk,
        font=\PortraitTitlePageMediumLatin\PortraitTitlePageMediumCJK\fontsize{12}{15}\selectfont]
        at (0,55mm) {\@author};
      \ifx\PortraitTitlePage@Image\@empty\else
        \node[inner sep=0pt] at (0,-20mm) {\includegraphics[
          width=\PortraitTitlePageImageWidth,height=\PortraitTitlePageImageHeight,
          keepaspectratio]{\PortraitTitlePage@Image}};
      \fi
      % 底边对齐：提交说明换行时向上增长，不进入下边距。
      % flush center 使用 fil 胶居中；align=center 的有限伸缩会拉开 xeCJK 字距。
      \PortraitTitlePageIfCommitTF{%
        \node[anchor=south,inner sep=0pt,text width=150mm,align=flush center,
          text=PortraitTitlePageSecondary,font=\sffamily\fontsize{8}{11}\selectfont]
          at (0,-131.5mm) {\PortraitTitlePageCommitText};
      }{}%
    \end{scope}
  \end{tikzpicture}%
  \clearpage}
\makeatother
